%%%%%%%%%%%%%%%%%%%%%%%%%%%%%%%%%%%%%%%%%%%%%%%%%%%%%%%%%%%%%%%%%%%%%%%%%%%%%%%
%% Lista de Siglas e Abreviaturas (pacote glossaries)
%%
%% Contém apenas as siglas efetivamente usadas nos Capítulos 1 a 5 e nos
%% apêndices. Acrescentar entrada aqui ao introduzir sigla nova no texto.
%%%%%%%%%%%%%%%%%%%%%%%%%%%%%%%%%%%%%%%%%%%%%%%%%%%%%%%%%%%%%%%%%%%%%%%%%%%%%%%

%% Instrumentação e mecânica
\newacronym{MEMS}{MEMS}{sistemas microeletromecânicos (\emph{Micro-Electro-Mechanical Systems})}
\newacronym{IMU}{IMU}{unidade de medição inercial (\emph{Inertial Measurement Unit})}
\newacronym{BPFO}{BPFO}{frequência de passagem de esferas na pista externa (\emph{Ball Pass Frequency, Outer race})}
\newacronym{BPFI}{BPFI}{frequência de passagem de esferas na pista interna (\emph{Ball Pass Frequency, Inner race})}
\newacronym{DLPF}{DLPF}{filtro passa-baixas digital (\emph{Digital Low-Pass Filter})}
\newacronym{RMS}{RMS}{valor eficaz (\emph{Root Mean Square})}
\newacronym{ASA}{ASA}{acrilonitrila estireno acrilato}
\newacronym{NTC}{NTC}{termistor de coeficiente de temperatura negativo}

%% Normalização
\newacronym{ISO}{ISO}{Organização Internacional de Normalização (\emph{International Organization for Standardization})}
\newacronym{ABNT}{ABNT}{Associação Brasileira de Normas Técnicas}

%% Detecção e avaliação
\newacronym{ROC}{ROC}{característica de operação do receptor (\emph{Receiver Operating Characteristic})}
\newacronym{AUC}{AUC}{área sob a curva (\emph{Area Under the Curve})}
\newacronym{pAUC}{pAUC}{área parcial sob a curva, restrita a uma faixa de falso alarme}
\newacronym{PRAUC}{PR-AUC}{área sob a curva de precisão e revocação}
\newacronym{TPR}{TPR}{taxa de detecção (\emph{True Positive Rate})}
\newacronym{DFT}{DFT}{transformada discreta de Fourier (\emph{Discrete Fourier Transform})}
\newacronym{FFT}{FFT}{transformada rápida de Fourier (\emph{Fast Fourier Transform})}
\newacronym{FPR}{FPR}{taxa de falso alarme (\emph{False Positive Rate})}

%% Computação embarcada e comunicação
\newacronym{TinyML}{TinyML}{aprendizado de máquina em dispositivos de recursos limitados}
\newacronym{LoRa}{LoRa}{modulação de longo alcance e baixa potência (\emph{Long Range})}
\newacronym{SF}{SF}{fator de espalhamento (\emph{Spreading Factor})}
\newacronym{CRC}{CRC}{verificação cíclica de redundância (\emph{Cyclic Redundancy Check})}
\newacronym{ADC}{ADC}{conversor analógico-digital}
\newacronym{SRAM}{SRAM}{memória estática de acesso aleatório (\emph{Static Random-Access Memory})}

%% Projeto de hardware
\newacronym{PCB}{PCB}{placa de circuito impresso (\emph{Printed Circuit Board})}
\newacronym{ERC}{ERC}{verificação de regras elétricas (\emph{Electrical Rules Check})}
\newacronym{DRC}{DRC}{verificação de regras de projeto (\emph{Design Rules Check})}
\newacronym{PDN}{PDN}{rede de distribuição de alimentação (\emph{Power Delivery Network})}

%% Instituição
\newacronym{UFRJ}{UFRJ}{Universidade Federal do Rio de Janeiro}

%%%%%%%%%%%%%%%%%%%%%%%%%%%%%%%%%%%%%%%%%%%%%%%%%%%%%%%%%%%%%%%%%%%%%%%%%%%%%%%
%% O pacote glossaries só imprime as siglas EFETIVAMENTE USADAS via \gls{...}.
%% Como o texto dos capítulos usa as siglas em forma literal (e não \gls{}),
%% \glsaddall força todas as entradas acima para a lista impressa.
%%
%% Se no futuro o texto passar a usar \gls{}, remova a linha abaixo para que a
%% lista volte a conter apenas o que é citado.
%%%%%%%%%%%%%%%%%%%%%%%%%%%%%%%%%%%%%%%%%%%%%%%%%%%%%%%%%%%%%%%%%%%%%%%%%%%%%%%
\glsaddall
